\documentclass[12pt, a4paper]{article}
\usepackage{amsmath}
\usepackage{amssymb}
\usepackage{graphicx}
\usepackage{geometry}
\usepackage{listings}
\usepackage{xcolor}

\geometry{margin=1in}

\lstset{
    language=Matlab,
    basicstyle=\ttfamily\small,
    keywordstyle=\color{blue},
    stringstyle=\color{purple},
    commentstyle=\color{teal},
    numbers=left,
    numberstyle=\tiny\color{gray},
    stepnumber=1,
    breaklines=true,
    frame=single,
    backgroundcolor=\color{lightgray!10}
}

\title{Computer Assignment for Chapter 6 \\ \large Monte Carlo Estimation of Option Pricing in a Binomial Model}
\author{Matteo Pinzani}
\date{March 2026}

\begin{document}

\maketitle

\section{Introduction}
This report presents the estimation of a European Call option price using the Monte Carlo method within a discrete-time binomial model. We are given a two-step model ($N = 2$) and use random data to first calculate the exact theoretical price via the risk-neutral valutation formula, and subsequently approximate it through simulation.

\section{Theoretical Background and Exact Pricing}
The market consists of a risk-free bond and a risky asset (stock). The parameters are:
\begin{itemize}
    \item Initial stock price: $S(0) = 120$
    \item Upward return factor: $U = 0.2$
    \item Downward return factor: $D = -0.1$
    \item Risk-free rate: $R = 0.1$
    \item Strike price: $X = 120$
    \item Number of steps: $N = 2$
\end{itemize}

\subsection{Risk-Neutral Probabilities}
To avoid arbitrage ($D < R < U$), the pricing of contingent claim is performed under the risk-neutral measure $\mathbb{P}^*$. The risk-neutral probability of an upward movement, $p^*$, is given by:
\begin{equation}
    p^* = \frac{R - D}{U - D} = \frac{2}{3} \approx 0.6667
\end{equation}
Consequently, the probability of a downward movement is $q^* = 1 - p^* = \frac{1}{3}$.

\subsection{Final Stock Prices and Returns}
At the expity time $T = 2$, the stock can take three possible value depending on the path:
\begin{align*}
    S_{uu}(2) &= S(0)(1+U)^2 = 172.8 \\
    S_{ud}(2) &= S(0)(1+U)(1+D) = 129.6 \\
    S_{dd}(2) &= S(0)(1+D)^2 = 97.2
\end{align*}
The return of the European Call option is $C(2) = \max(S(2) - X, 0)$, producing:
\begin{align*}
    C_{uu}(2) &= \max(S_{uu}(2) - X, 0) =  52.8 \\
    C_{ud}(2) &= \max(S_{ud}(2) - X, 0) = 9.6 \\
    C_{dd}(2) &= \max(S_{dd}(2) - X, 0) = 0
\end{align*}

\subsection{Discounted Expectation}
The valuation relies on the fundamental no-arbitrage theorem for discrete-time models. According to this principle, the cost of a self-financing strategy that replicates a contingent claim must equal the fair price of the claim itself. Consequently, the inital price of the option, $C_E(0)$, can be expressed simply as the discounted expected value of its final return, calculated under the risk-neutral probability measure $\mathbb{P}^*$:
\begin{equation}
    C_E(0) = \frac{1}{(1+R)^2} \mathbb{E}^*[C(2)]
\end{equation}
Expanding the expectation:
\begin{equation}
    \mathbb{E}^*[C(2)] = (p^*)^2 C_{uu}(2) + 2p^*q^* C_{ud}(2) + (q^*)^2 C_{dd}(2) = 27.7333 \\
\end{equation}
Discounting back to $t = 0$:
\begin{equation}
    C_E(0) = \frac{\mathbb{E}^*[C(2)]}{(1+R)^2} = 22.9201
\end{equation}

\section{Monte Carlo Estimation}
The Monte Carlo methos estimates the expected return $\mathbb{E}^*[C(2)]$ by simulating a large number of independent random paths for the stock price using the risk-neutral probabilities.\\
Let $M$ be the total number of simulated paths. For the $i$-th simulation, the final stock price $S_i(2)$ is generated, and the corresponding return $f(S_i(2)) = \max(S_i(2) - X, 0)$ is calculated.\\
By the Law of Large Numbers, the expectation is approximated by the sample mean:
\begin{equation}
    C_E(0) \approx \frac{1}{(1+R)^2} \left[ \frac{1}{M} \sum_{i=1}^{M} f(S_i(2)) \right]
\end{equation}

\section{Results}
Using MATLAB to simulate $M = 1,000,000$ paths, the algorithm randomly assigned an 'Up' or 'Down' movement at each step based on the uniform distribution threshold of $p^* = 2/3$.\\
The numerical results obtained are:
\begin{itemize}
    \item \textbf{Exact Analytical Price}: 22.9201
    \item \textbf{Monte Carlo Estimated Price}: $\approx$ 22.92
\end{itemize}
The Monte Carlo estimation converges tightly to the theoretical value, confirming the validity of both the risk-neutral pricing framework and the computational simulation.

\newpage
\section{MATLAB Source Code}
\begin{lstlisting}
% 1. Initialization of Data
S0 = 120;       % Initial stock price
U = 0.2;        % Upward return
D = -0.1;       % Downward return
R = 0.1;        % Risk-free rate
X = 120;        % Strike price
N = 2;          % Number of steps (T = 2)
M = 1000000;    % Number of Monte Carlo simulations

% 2. Calculate the Risk-Neutral Probabilities
p_star = (R - D) / (U - D);
q_star = 1 - p_star;    % or q_star = (U - R) / (U - D)

% 3. EXACT Theoretical Calculation (for comparison)
% Final stock prices for the 3 possible nodes at T = 2
Suu = S0 * (1 + U)^2;
Sud = S0 * (1 + U) * (1 + D);
Sdd = S0 * (1 + D)^2;

% Corresponding returns
Cuu = max(Suu - X, 0);  % Cuu = (Suu - X)^+
Cud = max(Sud - X, 0);  % Cud = (Sud - X)^+
Cdd = max(Sdd - X, 0);  % Cdd = (Sdd - X)^+

% Expected payoff under risk-neutral measure E*[C(2)]
Expected_Return_Exact = (p_star^2)*Cuu + (2*p_star*q_star)*Cud + (q_star^2)*Cdd;

% Discounted expected return
C0_exact = Expected_Return_Exact / (1 + R)^2;

% 4. Monte Carlo Simulation
% Simulate step 1 and step 2 for M paths using uniformly distributed random
% numbers
step1_is_up = rand(M, 1) < p_star;
step2_is_up = rand(M, 1) < p_star;

% Total number of 'Up' steps for each simulated path (can be 0 -> [dd], 1 ->
% [ud, du], or 2 -> [uu])
up_steps = step1_is_up + step2_is_up;
down_steps = N - up_steps;

% Calculate the final stock prices for all M paths simultaneously
ST = S0 * ((1 + U).^up_steps) .* ((1 + D).^down_steps);

% Calculate returns for all M paths
Returns = max(ST - X, 0);

% Estimate the expected return E*[C(2)]
Expected_Return_MC = mean(Returns);

% Discount back to time 0
C0_MC = Expected_Return_MC / (1 + R)^2;

% 5. Display Results
fprintf('--- Binomial Option Pricing (N=2) ---\n');
fprintf('Risk-neutral probability p*:   %.4f\n', p_star);
fprintf('Exact Analytical Price:        %.4f\n', C0_exact);
fprintf('Monte Carlo Estimate:          %.4f\n', C0_MC);
fprintf('Estimation Error:              %.4f\n', abs(C0_exact - C0_MC));
\end{lstlisting}

\end{document}